\booksection{}{Préface courte}

Nous traversons un moment que beaucoup ressentent comme inhabituel, sans toujours parvenir à nommer exactement ce qui est en train de se produire. Depuis l'arrivée des modèles génératifs d'intelligence artificielle, l'idée qu'une machine puisse écrire, raisonner, programmer ou créer n'appartient plus à la science-fiction. Elle est devenue, en l'espace de quelques années seulement, une expérience ordinaire pour des millions de personnes.

Ce livre ne part pas du constat de ce que l'IA sait déjà faire. Il part d'une interrogation beaucoup plus simple : que se passe-t-il lorsque certaines capacités intellectuelles, jusqu'ici considérées comme rares ou réservées à des spécialistes, deviennent accessibles à une très large partie de l'humanité ?

Il faut d'emblée préciser un point essentiel. À l'heure où les discours sur l'IA oscillent entre un optimisme démesuré et un catastrophisme systématique, ce livre refuse l'un comme l'autre. Il n'a aucune prétention de prédire l'avenir. L'avenir de l'intelligence artificielle n'est pas écrit dans la technologie elle-même. Il dépendra, pour l'essentiel, des choix économiques, politiques, culturels et humains que nous serons collectivement capables de faire.

Ce n'est donc pas un livre sur l'IA. C'est un livre sur nous. Sur ce que nous valorisons, sur la manière dont nos sociétés ont historiquement organisé le travail, la connaissance, le pouvoir et sur les questions que nous devrons peut-être nous poser à nouveau si l'intelligence venait à devenir moins rare qu'elle ne l'a jamais été.
